% Solutions to the exercises of Chapter 23.
\begin{solutions}
  \item With 20 random points at 3--8~m and a random relative pose, $|\bar x_2\T E\bar x_1| \le 3\times10^{-17}$: rounding error. Moving one observation in the second image
        by 27~px gives an algebraic residual of $8\times10^{-3}$ and a Sampson distance $\sqrt{d_S}$ of $5.4$~px. Only the component of the error perpendicular to the
        epipolar line is visible (7.7~px of the 27~px here), and the Sampson distance spreads it over both images ($7.7/\sqrt2 = 5.4$). A two-view check cannot
        detect an error \emph{along} the epipolar line at all; a third view, or the filter's own gating, must catch those.

  \item Since $RR\T = I$ and $\sx{t}\T = -\sx{t}$, $EE\T = -\sx{t}^2 = \|t\|^2I - tt\T$ (\eqref{eq:skew-sq}). Every vector orthogonal to $t$ is an eigenvector with eigenvalue
        $\|t\|^2$, and $t$ itself has eigenvalue 0, so the singular values of $E$ are $(\|t\|, \|t\|, 0)$ and its rank is two. Degrees of freedom: $R$ has three and $t$ three,
        but the constraint $\bar x_2\T E\bar x_1 = 0$ is homogeneous, so $E$ and $\lambda E$ (i.e.\ $t$ and $\lambda t$) are indistinguishable: $3 + 2 = 5$. Scaling the scene and
        the baseline together leaves every bearing unchanged, so the scale of $t$ cannot be recovered from images. For the epipoles,
        \[
          t\T E = t\T\sx{t}R = -(\sx{t}t)\T R = -(t\times t)\T R = 0, \qquad E(R\T t) = \sx{t}t = 0 .
        \]
        The left null vector is $t$, the coordinates of the first camera center in the second frame ($X_1 = 0 \Rightarrow X_2 = t$), so $e_2 \propto t$ is its image. Likewise
        $e_1 \propto R\T t$ is the image of the second center ($X_2 = 0 \Rightarrow X_1 = -R\T t$).

  \item With 60 points, 0.5~px noise and 200 random scenes, recovering $(R, t)$ from $E = K\T FK$:
        \begin{center}\small
        \begin{tabular}{@{}l c c c@{}}
        \toprule
        & $\sigma_1/\sigma_8$ of $A$ & rotation error & direction of $t$ \\
        & & median / 90\% & median / 90\% \\
        \midrule
        raw pixels & $1.5\times10^5$ & $0.33^\circ$ / $1.6^\circ$ & $2.5^\circ$ / $11^\circ$ \\
        Hartley-normalized & 66 & $0.16^\circ$ / $0.34^\circ$ & $1.1^\circ$ / $2.4^\circ$ \\
        \bottomrule
        \end{tabular}
        \end{center}
        In pixels, the columns of $A$ range from $1$ to $10^5$, so the smallest singular vector is dominated by rounding and by the noise in the large entries.
        Normalization makes all columns comparable. It halves the typical error and removes the bad cases, which matter most inside RANSAC.

  \item Take the SVD of $E$, flip the sign of $U$ or $V$ if its determinant is $-1$ (allowed, since $E$ is defined up to sign), and form the four candidates. For each,
        triangulate every correspondence (two depths from $\lambda_2\bar x_2 = \lambda_1R\bar x_1 + t$ by least squares) and count the points with both depths positive. In 500 random
        scenes, the true solution won every time, with all points in front of both cameras. For each wrong candidate, \emph{all} points failed: in a typical scene one
        candidate put them in front of the first camera but behind the second, and the other two put them behind the first. With noisy data, use the count, not
        a requirement that every point pass.

  \item With $R$ known, each correspondence gives $c_i = R\bar x_1^{(i)}\times\bar x_2^{(i)}$, the hypothesis is $t \propto c_i\times c_j$, and inliers have $\sqrt{d_S}$ below the threshold (2~px).
        For 100 correspondences with 1-pixel noise and 30\% outliers, adaptive RANSAC stopped after 9.8 iterations on average ($\lceil\log0.01/\log(1 - 0.7^2)\rceil = 7$
        at the true ratio; the observed ratio is slightly lower because the threshold rejects some inliers). After refitting $t$ to all inliers (the unit $t$
        minimizing $\sum(t\T c_i)^2$), the median direction error was $0.6^\circ$; 4.6\% of the inliers were rejected (the expected rate of a $2\sigma$ threshold) and 1.3\% of the
        outliers accepted (those lying near their epipolar line by chance). The five-point method needs $\lceil\log0.01/\log(1 - 0.7^5)\rceil = 26$ iterations at the same confidence,
        each solving a degree-10 polynomial and scoring up to ten candidates, against one cross product. The price of the two-point method is its dependence
        on the gyroscope: a rotation error $\delta$ shifts predicted points by about $f\delta$ pixels, so the threshold must cover the gyro error accumulated between
        frames.

  \item With $t = 0$, every correspondence satisfies $\bar x_2 \propto R\bar x_1$, so $\bar x_2\T\sx{t}R\bar x_1 = \bar x_2\T(t\times R\bar x_1) = 0$ for \emph{every} $t$. The data fit the
        whole family $E = \sx{t}R$, and the eight-point matrix $A$ has three singular values at the noise level instead of one. In simulation, a randomly chosen
        $t$ combined with the true $R$ fits 68--70\% of the tracks to within 1~px, exactly the fraction the true model fits with 1-pixel noise. Whatever
        \code{cv::findEssentialMat} returns therefore has an arbitrary translation direction that changes from run to run, and its inlier set is no test of the
        tracks, since the free $t$ can turn epipolar lines through outliers. The parallax test: the median of $\|\pi(R\bar x_1) - \pi(\bar x_2)\|$ over the tracks was
        $1.6$~px under pure rotation (noise only) and $8.8$~px for a 10~cm sideways motion with points at 3--8~m. A threshold of a few pixels separates the two
        regimes. Below it, check tracks only against $\bar x_2 \propto R\bar x_1$.

  \item The matrix $\sx{\bar x_i}$ has rank two (its nullspace is $\bar x_i$), and multiplying by the invertible $R_{wc_i}\T$ keeps the rank, so each view contributes three
        equations of rank two. For two views, $p$ is in the nullspace of the stacked matrix if and only if $R_{wc_i}\T p \parallel \bar x_i$ for both views, i.e.\ $p$ is parallel to both ray
        directions $r_i = R_{wc_i}\bar x_i$. If the rays are not parallel, only $p = 0$ qualifies: the $6\times3$ matrix has full column rank, and the least-squares solution
        is unique. Parallel rays (zero parallax) leave a one-dimensional nullspace, the position along the ray.

  \item With the baseline perpendicular to the viewing direction, $Z = b/d$ with disparity $d = x_1 - x_2$. The two noises are independent, so $\sigma_d = \sqrt2\sigma_{uv}$, and to first
        order $\sigma_Z = |\partial Z/\partial d|\,\sigma_d = (b/d^2)\sqrt2\sigma_{uv} = (Z^2/b)\sqrt2\sigma_{uv}$. A Monte Carlo at $Z = 5$~m with 1-pixel noise reproduces the figure: relative errors
        of $17.7\%$, $8.8\%$, $3.5\%$, $1.8\%$ at parallax angles of $1^\circ$, $2^\circ$, $5^\circ$, $10^\circ$, matching the formula. For three views at $0$, $b/2$ and $b$, write the
        observations as $x_i = X/Z - s_i/Z + n_i$, with $s_i$ the camera positions along the baseline. The inverse depth $1/Z$ is the \emph{slope} of a straight-line fit of $x_i$
        against $s_i$, with variance $\sigma_{uv}^2/\sum_i(s_i - \bar s)^2$. A view at the mean position $\bar s$ adds nothing to this sum: three views give the same depth accuracy as
        two (simulation: within 2\%), and five equally spaced views improve it by only 11\%. Depth accuracy is bought with baseline; intermediate views
        add redundancy, which is what outlier detection and the MSCKF's nullspace projection need.

  \item Monocular bundle adjustment has seven gauge freedoms: a global translation (3), a global rotation (3) and a global scale (1) leave all reprojection errors
        unchanged. Visual--inertial bundle adjustment has four: translation (3) and rotation about the gravity direction (1). The accelerometer measures
        specific force in the body frame, and the residuals compare it with $R_{wb}\T(a_w - g_w)$ for the fixed $g_w$. A global rotation about a horizontal axis would tilt the
        trajectory relative to gravity and change these residuals, so roll and pitch are observable. A rotation about the vertical leaves $g_w$ invariant and
        changes nothing. Scale becomes observable because the IMU measures metric accelerations, which a scaled trajectory would contradict, provided
        the motion actually accelerates (at constant velocity, scale is unobservable again). Consequently, fix the gauge of a visual--inertial problem with
        the first position and yaw only: fixing roll and pitch as well would inject false information.
\end{solutions}
